\documentclass[11pt]{article}
\usepackage[T1]{fontenc}
\usepackage{lmodern,amsmath,amssymb,amsthm,mathtools,microtype}
\usepackage[margin=1in]{geometry}
\usepackage{enumitem}
\usepackage[colorlinks=true,linkcolor=blue,citecolor=blue,urlcolor=blue]{hyperref}
\newtheorem{theorem}{Theorem}
\newtheorem{lemma}[theorem]{Lemma}
\newtheorem{corollary}[theorem]{Corollary}
\newtheorem{proposition}[theorem]{Proposition}
\theoremstyle{definition}
\newtheorem{definition}[theorem]{Definition}
\newcommand{\Fp}{\mathbb F_p}
\newcommand{\Mat}{\operatorname{M}}
\newcommand{\diag}{\operatorname{diag}}
\newcommand{\Aug}{\Delta}
\newcommand{\GL}{\operatorname{GL}}
\newcommand{\E}{\operatorname{E}}
\newcommand{\Span}{\operatorname{span}}
\newcommand{\Aut}{\operatorname{Aut}}
\title{Locally finite \(p\)-groups in Golod groups\\and an AT section}
\author{Alec Kriebel\thanks{ORCID: \href{https://orcid.org/0009-0001-9320-500X}{0009-0001-9320-500X}.}}
\date{24 October 2026 \quad Version 1.0.0}
\begin{document}
\maketitle
\begin{abstract}
We give an affirmative answer to Kourovka Notebook Problem 21.131 of
A.~V.~Timofeenko. For every prime \(p\), we construct a countable infinite
locally finite AT \(p\)-group satisfying Rozhkov's original infinite-support
definition. More generally, every countably infinite locally finite
\(p\)-group embeds in a Golod group in the precise nil-algebra sense of
Problem 9.76. It is also a quotient of a subgroup of a classical homogeneous,
residually finite Golod group. The argument combines augmentation ideals
with the stable-nil embedding construction of
Alahmadi--Alsulami--Jain--Zelmanov. Relative elementary matrices place the
desired derived subgroup inside a finitely generated canonical adjoint
subgroup. Tracking scalar extension permits Hammoudi's homogeneous Golod
cover to be applied.
\end{abstract}

\section{The problem and the precise source groups}

Problem 21.131 in the twenty-first edition of the Kourovka Notebook
\cite[p.~199]{KN}, proposed by A.~V.~Timofeenko, asks for an epimorphism
from a subgroup of a Golod group, as defined in Problem 9.76, onto an
infinite AT-group in the sense of Rozhkov \cite{Rozhkov}. We answer this
existential question. The AT target below is infinitely generated.
The whole source Golod group is not asserted to be AT, so the construction
does not settle the distinct Problem 13.55.

All algebras are associative and may be nonunital. For a nil algebra \(C\),
write \(C^\#=\Fp\oplus C\) for its unitization. Its adjoint group is \(1+C\).
For finite algebra generators \(c_1,\ldots,c_r\), put
\[
 G(C;c_1,\ldots,c_r)=\langle 1+c_1,\ldots,1+c_r\rangle\leq 1+C.
\]
If \(r\geq2\) and \(C\) is nil but not nilpotent, this is exactly a
Golod group in the definition of \cite[Problem 9.76, p.~38]{KN}: present
\(C\) as a quotient of the free nonunital algebra on these generators.
That definition does not require a homogeneous ideal or residual
finiteness. We shall additionally obtain the requested section using
a homogeneous classical Golod source.

Call an algebra \emph{stable nil} if \(\Mat_n(C)\) is nil for every finite
\(n\). In this paper, \emph{absolutely stable nil} means that
\(\Mat_n(C\otimes_F K)\) is nil for every field extension \(K/F\) and
every finite \(n\). In particular it is absolutely nil.
These last two properties are not interchanged without proof.

\begin{theorem}\label{thm:main}
Let \(p\) be a prime and \(H\) a countably infinite locally finite
\(p\)-group.
\begin{enumerate}[label=(\roman*),nosep]
\item \(H\) embeds in a Golod group in the precise sense of Problem 9.76.
\item There are a finitely generated, residually finite classical Golod
\(p\)-group \(\Gamma\), a subgroup \(P\leq\Gamma\), and an epimorphism
\(P\twoheadrightarrow H\). The defining nil algebra of \(\Gamma\) may be
chosen positively graded with homogeneous Golod relations.
\end{enumerate}
\end{theorem}

\begin{corollary}\label{cor:AT}
For every prime \(p\), a subgroup of a residually finite classical Golod
\(p\)-group maps onto an infinite AT-group in Rozhkov's original definition.
Consequently Problem 21.131 has an affirmative answer.
\end{corollary}

Phillips \cite{Phillips} already established general embeddings of
countable periodic groups into finitely generated periodic groups.
The issue addressed here is the canonical nil-adjoint realization.
The algebraic embedding input is due to
Alahmadi--Alsulami--Jain--Zelmanov \cite{AAJZshort,AAJZ};
the passage to a homogeneous cover uses Hammoudi \cite{Hammoudi}.
We make these dependencies explicit.

\section{A locally finite AT target}

Use the rooted \(p\)-ary tree with alphabet
\(A=\{0,1,\ldots,p-1\}\). A portrait assigns a permutation of \(A\)
at every vertex. Let \(a\) have the cycle
\(\sigma=(0\,1\,\cdots\,p-1)\) at the root and identity elsewhere.
Number tuple coordinates from zero and tree depths from zero.

We recall precisely the part of \cite[\S1]{Rozhkov} used here.
A rooted mutation acts only on coordinate zero. A longitudinal mutation
fixes an infinite directing ray, has infinite support in coordinate
levels, and, on a tuple first departing from that ray at coordinate \(j\),
changes only coordinate \(j+1\), by a permutation depending on the mutation,
the level and the departing symbol. An AT-group is generated by a subgroup
\(C\) of rooted mutations and a subset \(D\) of longitudinal mutations,
with all accompanying permutation groups transitive on the respective
alphabets. In particular \(D\) need not be finite or be a subgroup.

Let \(i(1),i(2),\ldots\) be the concatenation
\[
 1;\quad 1,2;\quad 1,2,3;\quad 1,2,3,4;\quad\cdots.
\]
For \(i\geq1\), define \(d_i\) to have portrait \(\sigma\) at precisely
the vertices
\[
 v_n=0^{n-1}1,\qquad i(n)=i,
\]
and identity at all other vertices. The vertices \(v_n\) are pairwise
incomparable. Hence all these local actions have disjoint supports.
The \(d_i\) commute, have order \(p\), and are independent. Set
\[
 H=\langle a,d_1,d_2,\ldots\rangle.
\]

Every \(d_i\) fixes \(0^\infty\). On a tuple whose first nonzero coordinate
is \(j\), it changes only coordinate \(j+1\): it acts by \(\sigma\)
if the departing symbol is \(1\) and \(i(j+1)=i\), and otherwise acts
trivially. Each \(i\) occurs infinitely often, so every designated \(d_i\)
has infinite level support. Take \(C=\langle a\rangle\) and
\(D=\{d_i:i\geq1\}\). At level zero the accompanying group is
\(\langle\sigma\rangle\); at each level \(n\geq1\), the generator
\(d_{i(n)}\) supplies \(\sigma\) at \(v_n\). Thus every accompanying group
is \(\langle\sigma\rangle\), which is transitive. This proves the original
AT condition, including its infinite-support requirement.

\begin{lemma}\label{lem:target}
The group \(H\) is countable, infinite and locally finite \(p\).
\end{lemma}
\begin{proof}
Its independent subgroup \(\langle d_i:i\geq1\rangle\cong
\bigoplus_{i\geq1}C_p\) proves infinitude, and countability is immediate.
Fix a finite \(I\subset\mathbb N\) and
\(H_I=\langle a,d_i:i\in I\rangle\). Choose \(N\geq1\) with \(i(N)\notin I\).
Let \(t_i\) be the section of \(d_i\) at \(0^N\), and set
\(B_N=\langle t_i:i\in I\rangle\). Their remaining portraits are disjoint,
so \(B_N\) is finite elementary abelian, of order \(p^{|I|}\).

Every depth-\(N\) section of a chosen generator belongs to \(B_N\).
For \(a\) all such sections are trivial. For \(d_i\), the section at
\(0^N\) is \(t_i\). The only other possible nontrivial section is the
root cycle at \(v_N\); it is absent since \(i(N)\notin I\). Earlier
portraits have no descendants, and later portraits lie below \(0^N\).
The section formula
\((fg)|_v=f|_{g(v)}\,g|_v\), and its inverse version, now show that every
depth-\(N\) section of every element of \(H_I\) lies in \(B_N\).

The faithful decomposition into the action through depth \(N\) and its
depth-\(N\) sections embeds \(H_I\) in
\[
 B_N^{\,p^N}\rtimes W_N,
\]
where \(W_N\) is the finite iterated wreath product of \(N\) copies of
\(C_p\). This is a finite \(p\)-group. Every finite subset of \(H\)
is contained in one such \(H_I\), proving the assertion.
\end{proof}

The missing label at the chosen depth is essential: a selected label
there would leave a rooted cycle beside a directed tail. Local finiteness
is established on the infinite tree, not inferred from finite computations.

\section{From locally finite groups to adjoint derived subgroups}

\begin{lemma}\label{lem:envelope}
Every countable locally finite \(p\)-group \(H\) embeds in the derived
subgroup of a countable locally finite \(p\)-group \(K\).
\end{lemma}
\begin{proof}
Set \(K_0=H\), \(m_j=p^{j+1}\), and
\(K_{j+1}=K_j^{m_j}\rtimes\langle t_j\rangle\), where \(t_j\) has order
\(m_j\) and cyclically permutes the coordinates. Embed \(K_j\) diagonally
in the base and take the direct limit \(K\). Each stage, and hence the
limit, is countable and locally finite \(p\): finitely many coordinates
from finitely many elements are contained in one finite \(p\)-subgroup.

If \(h\) in a stage has order \(p^b\), use a later transition with
\(p^b\mid m_j\). Index its coordinates modulo \(m_j\) and arrange
\(t_j f t_j^{-1}(i)=f(i-1)\). For \(f(i)=h^i\) one has
\[
 (f t_j f^{-1}t_j^{-1})(i)=h^i h^{-(i-1)}=h.
\]
The wraparound equality uses \(h^{m_j}=1\).
Thus the image of every element of \(K\) lies in \(K'\), and \(K=K'\).
In particular the original \(H\) embeds in \(K'\).
\end{proof}

For the specific AT target one may use just \(K=H\wr C_p\):
each chosen generator has order \(p\), so the same computation puts its
diagonal image in \(K'\). It then puts the diagonal image of all \(H\)
in \(K'\), without asserting exponent \(p\) for \(H\).

\begin{lemma}\label{lem:augmentation}
For a countable locally finite \(p\)-group \(K\), the augmentation ideal
\(A=\Aug(\Fp[K])\) is countable dimensional and absolutely stable nil.
The map \(k\mapsto1+(k-1)\) embeds \(K\) in \(1+A\).
\end{lemma}
\begin{proof}
For any extension field \(E/\Fp\),
\(A\otimes E=\Aug(E[K])\). Finitely many elements of this ideal are
supported in a finite \(p\)-subgroup \(Q\).
Its augmentation ideal is nilpotent. Here is the standard induction:
choose a central \(z\in Q\) of order \(p\). The central ideal
\(J=(z-1)E[Q]\) satisfies \(J^p=0\), and
\(E[Q]/J\cong E[Q/\langle z\rangle]\).
If the augmentation ideal of the quotient has nilpotence index \(q\),
then \(\Aug(E[Q])^q\subseteq J\), so its \(pq\)-th power is zero.
This starts at the trivial group.

The entries of any finite matrix over \(\Aug(E[K])\) lie in one such
nilpotent augmentation ideal. The matrix is therefore nilpotent.
This proves absolute stable nilness. The group embedding is the usual
inclusion of the group basis into its group algebra and is injective.
\end{proof}

The first clause of \cite[Theorem 4.3]{AAJZ} embeds every countable-dimensional
stable nil algebra in a finitely generated stable nil algebra.
This result was already announced in \cite[Theorem 6]{AAJZshort}.
We first use this clause without a scalar-extension refinement.
Embedding \(A\) in such a \(B\) embeds \(K\) in \(U=1+B\);
Lemma~\ref{lem:envelope} therefore embeds \(H\) in \(U'\).
The extension to unitizations is injective, since it is the original
linear injection on the nil ideal and identity on the scalar summand.

\section{The finite canonical unit subgroup}

\begin{lemma}\label{lem:elementary}
Let \(B\) be a finitely generated stable nil \(\Fp\)-algebra.
There is a finitely generated group \(L\leq1+\Mat_3(B)\) containing
an isomorphic copy of \((1+B)'\).
\end{lemma}
\begin{proof}
Put \(R=\Fp\oplus B\), and let \(s_1,\ldots,s_d\) generate \(B\) as a
nonunital algebra. Write \(e_{ij}(b)=I+bE_{ij}\). Let
\(E=\E_3(R)\), \(Q=\E_3(\Fp)\), and
\(X=\{e_{ij}(s_\ell):i\ne j,\ 1\leq\ell\leq d\}\).
For distinct indices the identities
\[
 e_{ij}(b)e_{ij}(c)=e_{ij}(b+c),\qquad
 [e_{ik}(b),e_{kj}(c)]=e_{ij}(bc)
\]
hold with commutator \(xyx^{-1}y^{-1}\).
Induction on words, followed by addition and integer powers, shows that
\(\langle X\rangle\) contains every \(e_{ij}(b)\), \(b\in B\).
The use of the prime field ensures all scalar coefficients are obtained
by integer powers. Since
\(e_{ij}(\lambda+b)=e_{ij}(\lambda)e_{ij}(b)\), \(E=\langle Q,X\rangle\).

Define
\[
 L=\langle qxq^{-1}:q\in Q,\ x\in X\rangle.
\]
It is finitely generated because \(Q\) is finite. It contains \(X\),
is normalized by \(Q\), and is trivially normalized by itself.
Hence it is normal in \(E=\langle Q,L\rangle\); it equals the normal
closure of the elementary matrices over \(B\). Each displayed generator
lies in \(1+\Mat_3(B)\), as does \(L\).

For \(u\in1+B\), put \(D(u)=\diag(u,u^{-1},1)\). Elementary
factorization gives
\[
 D(u)=e_{12}(u)e_{21}(-u^{-1})e_{12}(u)
             e_{12}(-1)e_{21}(1)e_{12}(-1).
\]
In \(E/L\), each factor may be replaced by its reduction modulo \(B\).
The resulting product is its value at \(u=1\), namely identity.
Thus \(D(u)\in L\). Direct multiplication, retaining the coefficient order,
gives
\[
 D(u)D(v)D(uv)^{-1}=\diag(1,u^{-1}v^{-1}uv,1).
\]
These coefficients generate \((1+B)'\). Therefore
\(h\mapsto\diag(1,h,1)\) embeds \((1+B)'\) in \(L\).
No homomorphism assertion for \(u\mapsto D(u)\) is used.
\end{proof}

\begin{proof}[Proof of Theorem~\ref{thm:main}(i)]
Use the preceding lemmas to embed \(H\) in \(L\).
List its finite displayed generators as \(I+c_1,\ldots,I+c_r\) and let
\(C\) be the nonunital subalgebra of \(\Mat_3(B)\) generated by the \(c_i\).
It is nil because \(\Mat_3(B)\) is nil.
Every inverse of \(I+c_i\) is a finite geometric polynomial in \(c_i\);
hence \(L=G(C;c_1,\ldots,c_r)\leq1+C\).
If \(C\) were nilpotent, words below a fixed length would span it.
It would be finite dimensional over the finite field, making \(1+C\)
finite, contrary to \(H\leq L\).

All elements \(1+c\) have finite \(p\)-power order:
\((1+c)^{p^e}=1+c^{p^e}=1\) for large \(e\).
Thus infinite \(L\) is not cyclic. After deleting identity entries and
repetitions its finite list has at least two entries; the remaining
differences still generate the same algebra. This gives precisely the
source required in Problem 9.76.
\end{proof}

\section{A homogeneous residually finite source}

For clarity we verify the scalar-extension property needed for Hammoudi.
Stable nilness alone is not substituted for absolute nilness.

\begin{lemma}\label{lem:absolute}
A countable-dimensional absolutely stable nil \(\Fp\)-algebra \(A\)
embeds in a finitely generated absolutely stable nil \(\Fp\)-algebra.
\end{lemma}
\begin{proof}
We track the construction in \cite[Lemmas 4.1, 4.3--4.5 and proof of
Theorem 4.3]{AAJZ}. Its seed can be fixed before a field extension is chosen.
Indeed their Lemma 4.3, applied to
\[
 \Fp[t_1,t_2,\ldots]\otimes\Mat_\infty(\Fp),
\]
gives a fixed finitely generated infinite-dimensional algebra \(B_0\)
whose tensor product with this algebra is nil.
The polynomial-matrix algebra has locally subexponential growth:
a finite set lies in a finite matrix corner over finitely many polynomial
variables, where growth is polynomial.
For every \(E/\Fp\) and every finite sum
\(X=\sum_{\ell=1}^s b_\ell\otimes\alpha_\ell\otimes Y_\ell
\in\Mat_n(B_0\otimes E)\), with \(b_\ell\in B_0\),
\(\alpha_\ell\in E\) and \(Y_\ell\in\Mat_n(\Fp)\), the universal element
\(\sum b_\ell\otimes t_\ell\otimes Y_\ell\) is nilpotent.
Evaluation \(t_\ell\mapsto\alpha_\ell\) proves \(X\) nilpotent.
Thus this same \(B_0\) is absolutely stable nil. This is the fixed-seed
argument in their Lemma 4.4.

Their auxiliary algebra is
\(\widetilde A=A\wr N\), where \(N=\Fp u+\Fp u^2\), \(u^3=0\).
It is the external semidirect sum \(N+\Mat_2(A)\).
Base change commutes with this finite-dimensional wreath construction.
For every \(E\) and \(n\), its matrix ideal is
\(\Mat_{2n}(A\otimes E)\), nil by assumption, and its quotient is
\(\Mat_n(N\otimes E)\), nilpotent. An extension of a nil ideal by a nil
algebra is nil: a power of each element enters the ideal, and another
power is zero. Hence \(\widetilde A\) is absolutely stable nil.

Put \(V=\Fp\oplus B_0\), with a fixed basis, and let
\(S_F=S(\widetilde A,V)\) be the matrices over \(\widetilde A\)
with finitely many nonzero rows, as in \cite{AAJZ}.
The natural map
\[
 S_F\otimes E\longrightarrow S(\widetilde A\otimes E,V\otimes E)
\]
is injective and multiplicative. To check injectivity, write a tensor as
\(\sum M_\ell\otimes\alpha_\ell\) with the \(\alpha_\ell\)
\(\Fp\)-linearly independent; zero in each matrix entry then forces
all entries of each \(M_\ell\) to vanish. Finite sums of matrices have a
common finite row support. Matrix multiplication also involves finite
sums, because the right factor has finite row support.
No surjectivity of this scalar-extension map is required.

The bimodule actions are left multiplication on \(V\) and precomposition
by left multiplication on \(V\); these extend to the same actions over
\(E\). Consequently the external semidirect sums give an injection
\[
 (B_0+S_F)\otimes E\ \longrightarrow\
 R_E=(B_0\otimes E)+S(\widetilde A\otimes E,V\otimes E).
\]
The finite-row algebra is an ideal: postcomposition sends a finite-dimensional
output space into another such space, while precomposition preserves its
output space.

Every \(\Mat_n\) of that finite-row ideal is nil. View its blocks as
one finite-row matrix indexed by \(\{1,\ldots,n\}\) times the chosen basis.
For a finite row support \(T\), the \(T\times T\) corner map is a
homomorphism whose kernel has square zero. The corner is nil by stable
nilness of \(\widetilde A\otimes E\); therefore the original matrix is
nilpotent. The quotient \(\Mat_n(R_E)\) by this ideal is
\(\Mat_n(B_0\otimes E)\), also nil. Thus \(\Mat_n(R_E)\) is nil.

Finally the containment argument of \cite[proof of Theorem 4.3]{AAJZ}
places \(A\) in the finitely generated algebra
\[
 B_1=\langle B_0,c_\gamma,e_{11}(u)\rangle\leq B_0+S_F
\]
for one generating linear map \(\gamma:V\to\widetilde A\).
The two additional matrices have finite row support.
This invokes exactly their established embedding construction; only
its scalar-extension property is being checked here.
Field extension preserves the injection \(B_1\hookrightarrow B_0+S_F\).
The injection into \(R_E\) then proves every \(\Mat_n(B_1\otimes E)\)
nil. The finitely generated algebra \(B_1\) is the required one.
\end{proof}

\begin{proof}[Proof of Theorem~\ref{thm:main}(ii)]
Use Lemma~\ref{lem:absolute} for the augmentation ideal in
Lemma~\ref{lem:augmentation}, and use its absolutely stable nil containing
algebra as \(B\) in the elementary construction.
The resulting finitely generated difference algebra
\(C\subseteq\Mat_3(B)\) is absolutely nil: after every field extension it
injects into the nil algebra \(\Mat_3(B\otimes E)\).

Present \(C=F/I\), with \(F=\Fp\langle x_1,\ldots,x_r\rangle_+\)
and \(x_i\mapsto c_i\). Hammoudi's Theorem 3 \cite[p.~96]{Hammoudi}
provides a homogeneous ideal \(J\subseteq I\) such that \(D=F/J\)
is a classical Golod nil algebra. His construction retains these
generators. Therefore
\[
 \Gamma=\langle1+x_1+J,\ldots,1+x_r+J\rangle
   \ \xrightarrow{\ \rho\ }\ L
\]
is surjective; see also \cite[Corollary 3]{Hammoudi}.
The algebra \(D\) is positively graded, with finite-dimensional homogeneous
components. Its degree truncations are finite nilpotent \(\Fp\)-algebras
and separate its elements. The corresponding finite \(p\)-group maps
separate elements of \(\Gamma\), proving residual finiteness.
The group is infinite since it maps onto infinite \(L\).

For the embedding \(\iota:H\hookrightarrow L\) already constructed, set
\(P=\rho^{-1}(\iota(H))\). The map
\(\iota^{-1}\circ\rho|_P:P\twoheadrightarrow H\) is the required epimorphism.
\end{proof}

\begin{proof}[Proof of Corollary~\ref{cor:AT}]
Apply Theorem~\ref{thm:main}(ii) to the original AT group of
Lemma~\ref{lem:target}.
\end{proof}

\section{Scope, verification and disclosures}

The target and subgroup may be infinitely generated, as the problem
permits. We do not obtain a finitely generated AT quotient or an AT
realization of the entire source. The homogeneous-source conclusion
is a section statement; it does not assert that every countable locally
finite \(p\)-group embeds in a residually finite group.
For example a nontrivial finite image of the divisible Pr\"ufer
\(p\)-group is impossible.

The proofs above are theoretical. Supplementary standard-library scripts
check finite portraits, common-zero sections, nonabelian wreath identities,
augmentation powers and noncommutative elementary identities.
These finite checks do not prove the universal statements.
The package includes exact inputs, outputs and independent AI adversarial
reviews. External theorems remain identified as such; no formal
proof-assistant verification of this result is claimed.

The accompanying priority audit searched the official Notebook and
updates, arXiv, related primary literature and accessible bibliographic
records through 24 October 2026. It found no earlier complete same-scope
solution. This is a bounded assessment, not a certification of first
priority. The AT definition was checked using primary indexed full text;
repeated attempts to obtain the original 1986 PDF bytes failed.
The exact source excerpts and acquisition limitations are documented in
the local research record; third-party articles are not redistributed
with the verification package.

AI tools were used extensively in discovering, developing, checking and
writing the mathematics, including independent adversarial AI reviews.
This preprint has not undergone independent human peer review.

\begin{thebibliography}{9}
\bibitem{AAJZshort}
A.~Alahmadi, H.~Alsulami, S.~K.~Jain and E.~Zelmanov,
On matrix wreath products of algebras,
\emph{Electron. Res. Announc. Math. Sci.} \textbf{24} (2017), 78--86.
\href{https://doi.org/10.3934/era.2017.24.009}{doi:10.3934/era.2017.24.009}.

\bibitem{AAJZ}
A.~Alahmadi, H.~Alsulami, S.~K.~Jain and E.~Zelmanov,
Matrix wreath products of algebras and embedding theorems,
\emph{Trans. Amer. Math. Soc.} \textbf{372} (2019), 2389--2406.
\href{https://doi.org/10.1090/tran/7642}{doi:10.1090/tran/7642}.
Detailed theorem numbering used here is that of the author manuscript,
\href{https://arxiv.org/abs/1703.08734v2}{arXiv:1703.08734v2}.

\bibitem{Hammoudi}
L.~Hammoudi, Quotients of nilalgebras and their associated groups,
\emph{Pacific J. Math.} \textbf{212} (2003), no.~1, 93--101.
\href{https://doi.org/10.2140/pjm.2003.212.93}{doi:10.2140/pjm.2003.212.93}.

\bibitem{KN}
E.~I.~Khukhro and V.~D.~Mazurov (eds.),
\emph{Unsolved Problems in Group Theory: The Kourovka Notebook},
21st ed., Sobolev Institute of Mathematics, Novosibirsk, 2026,
official revision of 20 October 2026, Problems 9.76, 13.55 and 21.131.
\url{https://kourovkanotebookorg.wordpress.com/wp-content/uploads/2026/10/21tkt-3.pdf}.

\bibitem{Phillips}
R.~E.~Phillips, Embedding methods for periodic groups,
\emph{Proc. London Math. Soc.} (3) \textbf{35} (1977), 238--256;
correction, (3) \textbf{38} (1979), 152.
\href{https://doi.org/10.1112/plms/s3-35.2.238}{doi:10.1112/plms/s3-35.2.238}.

\bibitem{Rozhkov}
A.~V.~Rozhkov, Theory of Aleshin type groups,
\emph{Mat. Zametki} \textbf{40} (1986), no.~5, 572--589;
English translation, \emph{Math. Notes} \textbf{40} (1986), 827--836.
\href{https://doi.org/10.1007/BF01159699}{doi:10.1007/BF01159699}.
\end{thebibliography}
\end{document}
